% Chapter 08: Reference
% Complete field reference for manifest and binary interfaces

\chapter{Reference}
\label{ch:reference}

% ============================================================
\section{Kit Manifest Field Reference}
\label{sec:manifest-reference}
% ============================================================

This section provides a complete reference for all fields in the
\texttt{gert-kit.yaml} manifest.

\begin{longtable}{p{0.25\linewidth}p{0.12\linewidth}p{0.51\linewidth}}
\toprule
\textbf{Field} & \textbf{Required} & \textbf{Description} \\
\midrule
\endfirsthead

\toprule
\textbf{Field} & \textbf{Required} & \textbf{Description} \\
\midrule
\endhead

\midrule
\multicolumn{3}{r}{\footnotesize\itshape Continued on next page} \\
\bottomrule
\endfoot

\bottomrule
\endlastfoot

\texttt{apiVersion} & Yes & Must be \texttt{kit/v1}. \\

\texttt{meta.name} & Yes & Reverse-DNS namespaced identifier (e.g., \texttt{gert.compliance}, \texttt{com.acme.workflow}). \\

\texttt{meta.version} & Yes & Semantic version string (MUST follow semver, e.g., \texttt{"1.2.0"}). \\

\texttt{meta.description} & Yes & Human-readable summary of the Kit's purpose. \\

\texttt{meta.author} & No & Author name and contact (e.g., \texttt{"Acme Corp <kit@acme.example>"}). \\

\texttt{meta.homepage} & No & URL to Kit documentation or repository. \\

\texttt{meta.license} & No & SPDX license identifier (e.g., \texttt{"Apache-2.0"}). \\

\texttt{compatibility.gert-core} & No & Semver range specifying compatible gert core versions (e.g., \texttt{">=2.0.0 <3.0.0"}). If omitted, Kit claims compatibility with all versions. \\

\texttt{schema.apiVersion-suffix} & No & The suffix appended to \texttt{runbook/v1+} in runbooks authored with this Kit (e.g., \texttt{"compliance/v1"}). \\

\texttt{schema.json-schema} & No & Path (relative to Kit root) to the Kit's JSON Schema overlay (e.g., \texttt{"./schemas/compliance-v1.json"}). \\

\texttt{compiler.executable} & Yes & Path (relative to Kit root) to the compiler binary (e.g., \texttt{"./bin/gert-kit-compliance"}). \\

\texttt{compiler.args} & No & Array of arguments passed to the compiler (e.g., \texttt{["compile"]}). Default: no args. \\

\texttt{compiler.input-format} & No & Input format. Must be \texttt{yaml}. Default: \texttt{yaml}. \\

\texttt{compiler.output-format} & No & Output format. Must be \texttt{yaml}. Default: \texttt{yaml}. \\

\texttt{step-types} & No & Array of Kit-defined step types. Each entry has \texttt{name}, \texttt{lowers-to}, and \texttt{description}. \\

\texttt{step-types[].name} & Yes & Step type identifier (e.g., \texttt{approval-request}). \\

\texttt{step-types[].lowers-to} & Yes & Array of core step types this Kit type compiles into (e.g., \texttt{[manual, tool]}). \\

\texttt{step-types[].description} & No & Human-readable description of the step type. \\

\texttt{validators} & No & Array of domain-specific validators. \\

\texttt{validators[].name} & Yes & Validator identifier (e.g., \texttt{compliance/owner-required}). \\

\texttt{validators[].phase} & Yes & Validation phase: \texttt{parse} or \texttt{plan}. \\

\texttt{validators[].description} & No & Human-readable description of what the validator checks. \\

\texttt{projections} & No & Array of projections (read models over trace data). \\

\texttt{projections[].name} & Yes & Projection identifier (e.g., \texttt{compliance/audit-log}). \\

\texttt{projections[].description} & No & Human-readable description of the projection output. \\

\texttt{projections[].output-format} & No & Output format: \texttt{json}, \texttt{yaml}, \texttt{markdown}, or \texttt{text}. Default: \texttt{json}. \\

\end{longtable}

% ============================================================
\section{Kit Compiler Binary Interface}
\label{sec:compiler-reference}
% ============================================================

\subsection{Invocation}

\begin{verbatim}
<compiler-binary> [args...] < input.yaml > output.yaml
echo $?  # Exit code
\end{verbatim}

\subsection{Arguments}

Arguments are specified in the manifest under \texttt{compiler.args}. Common
convention:

\begin{description}
  \item[\texttt{compile}] --- Perform lowering (default command).
  \item[\texttt{validate}] --- Run validation only (optional, if validators are
    bundled into the compiler).
  \item[\texttt{--version}] --- Print Kit version and exit.
\end{description}

\subsection{Input}

Kit-authored runbook YAML on stdin. Example:

\begin{minted}{yaml}
apiVersion: runbook/v1+compliance/v1
meta:
  title: "Example runbook"
steps:
  - id: step-1
    type: approval-request
    approvers: [alice@acme.example]
    timeout: 1h
\end{minted}

\subsection{Output}

\paragraph{On success (exit code 0):}

Lowered \texttt{runbook/v1} YAML on stdout:

\begin{minted}{yaml}
apiVersion: runbook/v1
meta:
  title: "Example runbook"
  x-kit:lowered-from:
    kit: gert.compliance
    version: "1.2.0"
steps:
  - id: step-1
    type: manual
    approval:
      required: true
      approvers: [alice@acme.example]
      timeout: 3600
\end{minted}

\paragraph{On failure (exit code 1):}

JSON array of error messages on stdout:

\begin{minted}{json}
[
  {
    "file": "<stdin>",
    "line": 5,
    "severity": "error",
    "message": "approval-request requires at least one approver",
    "rule": "compliance/lowering"
  }
]
\end{minted}

\subsection{Exit Codes}

\begin{description}
  \item[\texttt{0}] --- Success. Stdout contains valid \texttt{runbook/v1} YAML.
  \item[\texttt{1}] --- Validation or lowering error. Stdout contains error JSON.
  \item[\texttt{2+}] --- Internal compiler error. Stderr may contain stack trace.
\end{description}

% ============================================================
\section{Validator Binary Interface}
\label{sec:validator-reference}
% ============================================================

\subsection{Invocation}

\begin{verbatim}
<validator-binary> < runbook.yaml > errors.json
echo $?  # Exit code: 0 = pass, 1 = errors
\end{verbatim}

\subsection{Input}

\paragraph{Parse-time validators:} Kit-authored runbook YAML on stdin.

\paragraph{Plan-time validators:} Lowered \texttt{runbook/v1} YAML on stdin.

\subsection{Output}

JSON array of validation messages:

\begin{minted}{json}
[
  {
    "file": "runbook.yaml",
    "line": 3,
    "column": 1,
    "severity": "error",
    "message": "Compliance runbooks must declare an owner in meta.roles.owner",
    "rule": "compliance/owner-required"
  },
  {
    "file": "runbook.yaml",
    "line": 15,
    "severity": "warning",
    "message": "Timeout exceeds recommended SLA",
    "rule": "compliance/timeout-sla"
  }
]
\end{minted}

If validation passes, output an empty array:

\begin{minted}{json}
[]
\end{minted}

\subsection{Exit Codes}

\begin{description}
  \item[\texttt{0}] --- No errors (warnings may be present).
  \item[\texttt{1}] --- At least one error.
  \item[\texttt{2+}] --- Internal validator error.
\end{description}

% ============================================================
\section{Projection Binary Interface}
\label{sec:projection-reference}
% ============================================================

\subsection{Invocation}

\begin{verbatim}
cat trace.ndjson | <projection-binary> > output.json
\end{verbatim}

\subsection{Input}

NDJSON trace event stream on stdin. Each line is a JSON object:

\begin{verbatim}
{"event":"run.started","timestamp":"2025-01-15T10:30:00Z","run_id":"r-abc123","payload":{}}
{"event":"step.started","timestamp":"2025-01-15T10:30:01Z","run_id":"r-abc123","step_id":"step-1","payload":{}}
{"event":"evidence.captured","timestamp":"2025-01-15T10:32:00Z","run_id":"r-abc123","step_id":"step-1","payload":{"type":"text","value":"..."}}
{"event":"run.completed","timestamp":"2025-01-15T10:35:00Z","run_id":"r-abc123","payload":{"status":"success"}}
\end{verbatim}

\subsection{Output}

Projection-specific format (JSON, YAML, Markdown, etc.) on stdout. Example
(audit log):

\begin{minted}{json}
{
  "run_id": "r-abc123",
  "audit_entries": [
    {
      "timestamp": "2025-01-15T10:32:00Z",
      "step_id": "step-1",
      "event_type": "evidence.captured",
      "details": {"type": "text", "value": "..."}
    }
  ]
}
\end{minted}

\subsection{Exit Codes}

\begin{description}
  \item[\texttt{0}] --- Success.
  \item[\texttt{1+}] --- Projection error (malformed input, internal error).
\end{description}

% ============================================================
\section{Environment Variables}
\label{sec:env-vars}
% ============================================================

\begin{description}
  \item[\texttt{GERT\_KIT\_PATH}] --- Colon-separated list of directories to
    search for Kits (similar to \texttt{PATH}). Example:
    \texttt{/usr/local/share/gert/kits:\textasciitilde/.gert/kits}.

  \item[\texttt{GERT\_KIT\_DEBUG}] --- If set to \texttt{1}, gert emits debug
    logs about Kit discovery and loading.

  \item[\texttt{GERT\_NO\_KIT\_CACHE}] --- If set to \texttt{1}, gert does not
    cache lowered runbooks. Useful during Kit development.
\end{description}

% ============================================================
\section{Error Codes and Meanings}
\label{sec:error-codes}
% ============================================================

\subsection{Compiler Exit Codes}

\begin{description}
  \item[\texttt{0}] --- Lowering succeeded.
  \item[\texttt{1}] --- Validation or lowering error (recoverable; error JSON
    on stdout).
  \item[\texttt{2}] --- Internal compiler error (panic, unexpected failure).
\end{description}

\subsection{Validator Exit Codes}

\begin{description}
  \item[\texttt{0}] --- Validation passed (no errors).
  \item[\texttt{1}] --- Validation failed (at least one error).
  \item[\texttt{2}] --- Internal validator error.
\end{description}

\subsection{Projection Exit Codes}

\begin{description}
  \item[\texttt{0}] --- Projection succeeded.
  \item[\texttt{1}] --- Malformed input (invalid NDJSON).
  \item[\texttt{2}] --- Internal projection error.
\end{description}

% ============================================================
\section{Trace Event Schema Reference}
\label{sec:trace-event-schema}
% ============================================================

All trace events have the following common fields:

\begin{description}
  \item[\texttt{event}] (string) --- Event type (e.g., \texttt{run.started},
    \texttt{step.completed}).
  \item[\texttt{timestamp}] (string) --- ISO 8601 timestamp (UTC).
  \item[\texttt{run\_id}] (string) --- Unique run identifier.
  \item[\texttt{step\_id}] (string, optional) --- Step ID (for step-level events).
  \item[\texttt{payload}] (object) --- Event-specific data.
\end{description}

\subsection{Event Types}

\begin{longtable}{p{0.28\linewidth}p{0.60\linewidth}}
\toprule
\textbf{Event Type} & \textbf{Payload Fields} \\
\midrule
\endfirsthead

\toprule
\textbf{Event Type} & \textbf{Payload Fields} \\
\midrule
\endhead

\midrule
\multicolumn{2}{r}{\footnotesize\itshape Continued on next page} \\
\bottomrule
\endfoot

\bottomrule
\endlastfoot

\texttt{run.started} & \texttt{runbook\_path}: string (path to runbook file) \\

\texttt{run.completed} & \texttt{status}: \texttt{success} | \texttt{failed} | \texttt{cancelled} \\

\texttt{step.started} & (empty) \\

\texttt{step.completed} & \texttt{exit\_code}: integer, \texttt{duration\_ms}: integer \\

\texttt{step.failed} & \texttt{error}: string (error message) \\

\texttt{evidence.captured} & \texttt{type}: \texttt{text} | \texttt{attachment}, \texttt{value}: string, \texttt{sha256}: string (for attachments) \\

\texttt{approval.requested} & \texttt{approvers}: array of strings \\

\texttt{approval.granted} & \texttt{approver}: string (who approved) \\

\texttt{approval.rejected} & \texttt{approver}: string, \texttt{reason}: string \\

\texttt{variable.set} & \texttt{name}: string, \texttt{value}: any \\

\end{longtable}

% ============================================================
\section{Common Pitfalls}
\label{sec:pitfalls}
% ============================================================

\paragraph{Pitfall 1: Non-deterministic lowering}

\textbf{Problem:} Compiler uses timestamps, random IDs, or environment variables
in lowered output.

\textbf{Solution:} Make the compiler pure. If you need unique IDs, derive them
deterministically from input data (e.g., hash of step ID + runbook title).

\paragraph{Pitfall 2: Validators with side effects}

\textbf{Problem:} Validator calls an API to check if a user exists.

\textbf{Solution:} Validators MUST be pure. If you need runtime checks, use a
\textbf{Tool} or \textbf{Extension}, not a validator.

\paragraph{Pitfall 3: Forgetting to handle validation phase}

\textbf{Problem:} Parse-time validator tries to validate cross-step invariants
(which requires seeing the full lowered runbook).

\textbf{Solution:} Use \texttt{phase: plan} for cross-step validators. They run
after lowering.

\paragraph{Pitfall 4: Breaking changes without major version bump}

\textbf{Problem:} Kit removes a field in version 1.3.0 (should be 2.0.0).

\textbf{Solution:} Follow semver strictly. Removing a field is a breaking
change → major version bump.

\paragraph{Pitfall 5: Not testing lowered output against core schema}

\textbf{Problem:} Compiler produces invalid \texttt{runbook/v1} YAML that gert
rejects.

\textbf{Solution:} Always validate lowered output against the core
\texttt{runbook/v1} JSON Schema as part of your test suite.
